\section{Choosing a Mode Map Without an Answer Key}
\label{sec:selector}

Everything so far was scored against ground truth. GraphCast offers none, and the
gap this opens is not a detail of evaluation but a hole in the middle of the
method. The pipeline of Sec.~\ref{sec:approach} begins by choosing a set of
variables---an aggregation map from the model's high-dimensional state to a
handful of candidate features---and that choice dominates everything downstream.
On SAVAR we can simply try many maps and keep the one with the best graph. On
GraphCast we cannot, so a recipe that needs the answer key to pick its own
variables is not a recipe at all.

This section builds a selector that ranks candidate maps using only quantities
computable without ground truth, and calibrates how much its verdict can be
trusted. It is the part of the project that has moved most since the results
above, and the part the GraphCast application depends on.

\subsection{The candidate pool and the two channels}
\label{sec:px}

The construction needs a pool of candidate aggregation maps spanning a wide
quality range, and two independent ways of reading a causal graph through each.

The pool (\texttt{sae/discover\_modes.py}) holds thirteen candidates: unsupervised
discovery operators (varimax-PCA, $k$-means and DMD clustering, each run on raw
pixels and on model activations), the oracle map, and deliberately corrupted
variants---blurred, shifted, merged, split, diagonal---that manufacture a
controlled quality axis. Scored against truth the pool spans F1 $0.00$ to $0.89$,
which is what a selector test requires.

Two channels then read a graph through each candidate map $\hat W$. The
\emph{int} channel pools the observation field through $\hat W$ and runs PCMCI+
over the resulting series, giving $\hat G_{\text{int}}$: this is discovery by
\emph{reading} correlations. The \emph{dyn} channel injects impulses into the
frozen forecaster along $\mathrm{pinv}(\hat W)$, reads the teacher-forced response
back out on the rows of $\hat W$, and thresholds against a pixel-permutation null,
giving $\hat G_{\text{dyn}}$: discovery by \emph{poking} the model. Neither uses
the true graph, the true modes, or the true footprints anywhere.

An earlier attempt at unsupervised selection scored candidates by the internal
consistency of a single channel (agreement between aggregate-level and
pixel-level independence verdicts). It reached Spearman $+0.52$ against true
quality and failed for a diagnosable reason: a map that destroys all signal has no
dependencies at either level and is therefore vacuously consistent. Consistency is
a trustworthy \emph{screen}---its zeros always mean a dead map---but not a
\emph{ranker}. That failure is what motivated crossing two channels instead.

\subsection{Same-map agreement fails, in two informative ways}

The natural two-channel statistic is the agreement between the two graphs read
through the \emph{same} map, $F_1(\hat G_{\text{int}}(A), \hat G_{\text{dyn}}(A))$.
It does not work: Spearman $+0.375$ against true quality, far below the
pre-registered bar of $0.8$. Table~\ref{tab:samew} shows why. The two failures are
structural, and both carry to GraphCast.

\begin{table}[h]
\centering
\small
\begin{tabular}{lccc}
\toprule
Candidate & True F1 & Same-$\hat W$ agreement & Pool-crossed (PX) \\
\midrule
blur      & 0.918 & 0.556 & 0.453 \\
oracle    & 0.879 & 0.621 & 0.447 \\
shift5    & 0.878 & \textbf{0.154} & 0.490 \\
vmax\_act & 0.835 & 0.298 & 0.424 \\
dmd\_act  & 0.632 & 0.118 & 0.351 \\
coarse4   & 0.230 & \textbf{0.727} & 0.149 \\
diag8     & 0.000 & 0.000 & 0.000 \\
\bottomrule
\end{tabular}
\caption{Selected candidates from the parent pool ($13$ total). Same-$\hat W$
agreement is maximized by the worst usable map (\texttt{coarse4}) and minimized by
one of the best (\texttt{shift5}); pool-crossing repairs both.}
\label{tab:samew}
\end{table}

First, \textbf{merged maps self-agree on the wrong graph}. The candidate
\texttt{coarse4}, which merges mode pairs and recovers almost nothing
(true F1 $0.230$), attains the pool's \emph{highest} same-map agreement
($0.727$). Both channels inherit the same marginalization from the shared $\hat W$
and therefore make the same mistake in the same place. Agreement measures
consistency of error, not correctness, whenever the error is shared.

Second, \textbf{the write channel is geometry-sensitive where the read channel is
not}. The candidate \texttt{shift5} has footprints in the wrong locations but
behaviourally coherent pooled series; it reads a nearly correct graph
(true F1 $0.878$) while its dyn channel returns almost nothing (agreement
$0.154$). Reading through $\hat W$ needs only that the pooled series be
behaviourally coherent; \emph{writing} through $\mathrm{pinv}(\hat W)$ needs the
injected perturbation to land on the true supports. This asymmetry is worth
stating on its own: perturbation-based validation implicitly certifies footprint
\emph{geometry}, not just behaviour, and a map can pass one test while failing the
other.

\subsection{Pool-crossed agreement}

Both failures come from the two channels sharing a map. Crossing the pool removes
the shared term. Define
\[
  \mathrm{PX}(A) \;=\; \frac{1}{|\mathcal{P}|-1}\sum_{B \neq A}
  F_1\!\left(\hat G_{\text{int}}(A),\, \hat G_{\text{dyn}}(B)\right),
\]
the mean agreement between candidate $A$'s read graph and every \emph{other}
candidate's poke graph. The dyn pool acts as a $\hat W$-independent reference, so
per-map channel error averages out while genuine map quality survives. The
direction matters: the int channel carries the quality signal, and the transposed
statistic (each candidate's dyn graph against the int pool) scores only $+0.288$.

On the parent rung PX reaches Spearman $+0.950$ (Pearson $+0.99$) against true
graph quality, against $+0.375$ for same-map agreement. Good maps sit at
$0.40$--$0.49$, merged or weak maps at $\le 0.17$, dead maps at exactly $0$---a
usable selector with a clean gap.

One caveat belongs here rather than in a footnote: the aggregation rule
(mean over the pool) was fixed \emph{after} the two confounds above were
diagnosed, on $13$ candidates. The parent number is therefore in-sample, and the
rungs below exist to test it out-of-sample.

\subsection{A scoring convention that had to be corrected}

Diagnosing \texttt{shift5} exposed an error in the ground truth itself, not in the
selector. Under the original convention---matching discovered variables to true
modes by footprint cosine---\texttt{shift5} scored F1 $0.000$. But its pooled
series correlate $0.979$ with the true modes; matching variables to modes by
series correlation instead gives F1 $0.835$. The map was good and the score was
wrong. The same correction lifts \texttt{dmd\_act} from $0.177$ to $0.553$.

All scoring in this section therefore uses behaviour-based matching. The
methodological point generalizes beyond SAVAR: \emph{footprint geometry is the
wrong success criterion for an aggregation map; behavioural coherence of the
pooled series is the right one}. On GraphCast, where features have no
ground-truth spatial support to be compared against, only the behavioural
criterion is available at all.

\subsection{Out-of-sample rungs}

Each rung in Table~\ref{tab:rungs} changes one property of the world or the
forecaster, retrains where needed, and reruns the battery unchanged. The rule was
pre-registered from the parent's saved battery and applied verbatim, with no
tuning; the bar is Spearman $\ge 0.8$.

\begin{table}[h]
\centering
\small
\begin{tabular}{llcc}
\toprule
Rung & What changes & PX Spearman & Verdict \\
\midrule
parent & --- (in-sample reference)            & $+0.950$ & --- \\
R1 & overlapping footprints (max cosine $0.20$) & $+0.946$ & \textbf{pass} \\
R3 & $C=2$ channels with cross-channel edges    & $+0.943$ & \textbf{pass} \\
R5 & rollout-trained forecaster (BPTT, horizons 1--8) & $+0.934$ & \textbf{pass} \\
R4b & $N=24$ modes, $N$ unknown a priori        & $+0.790$ & miss \\
\bottomrule
\end{tabular}
\caption{Pool-crossed selector across pre-registered rungs. Each is a full
generate--train--discover--calibrate chain; the dyn channel was verified live in
every one, so each number is a real test rather than an inapplicable channel.}
\label{tab:rungs}
\end{table}

The three passes matter for different reasons. \textbf{R1} removes the
disjoint-footprint assumption that every earlier SAVAR variant relied on, and it
is the sharpest test of the two-channel idea because overlap degrades the
internals channel specifically: discovery from activations collapses (\texttt{vmax\_act}
graph F1 $0.819 \to 0.185$) while the pixel side stays healthy ($0.938$). The
selector nonetheless ranks the pool correctly. \textbf{R5} fine-tunes the
forecaster on multi-step rollout, which is GraphCast's actual training curriculum;
it measurably moved the model's internals (\texttt{vmax\_act} truth-F1
$0.485 \to 0.556$) while leaving the pixel side identical, and the selector
survived the shift. \textbf{R3} adds coupled channels and cross-channel edges, and
passes even though its activation-side discovery is degenerate---the multivariate
GNN carries one hidden vector per \emph{spatial} node, so co-located channels
receive identical pooled activations. PX works there because it crosses the clean
pixel-side read channel against the live poke channel.

\textbf{R4b} is the honest miss. Tripling the mode count to $N=24$ with $N$
unknown, on a pre-registered resolution-homogeneous pool of ten candidates, PX
reaches $+0.790$ and misses its bar by $0.010$. It is reported as a miss and was
not re-run. Two things temper it. The Pearson correlation is $+0.975$, so the
calibration line is intact and only the ranking is noisy---with ten candidates a
single adjacent swap costs about $0.02$ Spearman. And the discovery side scaled
cleanly: unknown-$N$ community detection recovered $\hat N = 24$ and matched
$23/24$ modes at cosine $0.966$, while the fixed-rank varimax operators that
worked at $N=8$ collapsed entirely. That contrast is itself a result for the
GraphCast setting, where the number of features is large and not known in advance.

\subsection{Mechanism robustness}

A separate matrix (Table~\ref{tab:robust}) stresses the selector against the
realism mechanisms of the generator---non-Gaussian innovations, non-linear
dynamics, seasonality, and all of them combined---each as a full rung.

\begin{table}[h]
\centering
\small
\begin{tabular}{llcc}
\toprule
Mechanism & PX Spearman & Pearson & Verdict \\
\midrule
Non-Gaussian innovations (skew-normal) & $+0.909$ & $+0.946$ & \textbf{pass} \\
All mechanisms combined                & $+0.869$ & $+0.951$ & \textbf{pass} \\
Non-linearity (bilinear, $\alpha = 0.5$) & $+0.624$ & $+0.972$ & rank miss \\
Seasonality, raw series                & $+0.132$ & $+0.333$ & miss \\
Seasonality, deseasonalized            & $+0.769$ & $+0.789$ & near miss \\
\bottomrule
\end{tabular}
\caption{Selector under generator realism mechanisms. Same-$\hat W$ agreement
stays weak throughout ($+0.18$ to $+0.42$), so pool-crossing continues to carry
the signal.}
\label{tab:robust}
\end{table}

The combined world passing at $+0.869$ is the headline: with seasonality,
saturating and bilinear non-linearity, and skewed innovations all active at once,
the selector still ranks the pool. The two misses are more useful than the passes.

Under strong non-linearity PX becomes a screen rather than a ranker. It reliably
kills the dead maps, but eight candidates sit within $0.017$ of each other in PX
while spanning $0.29$ of true F1, and ranking inside that cluster is noise
(excluding two degenerate candidates the Spearman falls to $+0.382$). What
survives is the decision that actually matters: the top-pick costs only $0.057$ F1
relative to the best available map, and Pearson $+0.972$ says the calibration line
is undamaged. A related check confirms the linear conditional-independence test is
not the problem---under bilinear dynamics ParCorr is edge-recall-preserving, at a
cost of about four spurious edges, while a nonparametric test buys F1 $0.815 \to
0.957$ for a $50{,}000\times$ increase in compute.

Seasonality is the one mechanism that breaks the recipe outright on raw series
($+0.132$) and is repaired by preprocessing ($+0.769$). Shared periodic forcing
induces dependence between co-periodic modes that the conditional-independence
tests read as causal structure. \textbf{Deseasonalization is therefore a required
preprocessing step, not an option}---a concrete and slightly uncomfortable
finding for GraphCast, whose most prominent published SAE features include the
diurnal and seasonal cycles themselves~\cite{macmillan2025}.

\subsection{The trust dial and its per-domain calibration}
\label{sec:dial}

Ranking candidates is half of what is needed. The other half is turning an
observed PX into a predicted accuracy with an error bar---a dial that says how
much to believe a graph recovered from a model with no answer key.

Pooling all candidate points across the original five rungs ($n=52$) gives
\[
  \text{accuracy} \;\approx\; 1.48 \cdot \mathrm{PX} + 0.13,
  \qquad \text{Spearman } +0.905,\;\; R^2 = 0.797,\;\;
  \text{residual s.d.\ } 0.131 .
\]
The four mechanism worlds above were generated after this fit and never entered
it, so they constitute a genuine out-of-sample test ($n=52$ new points). The
result splits cleanly, and the split is the main finding of this subsection.

\emph{Ranking transfers}: Spearman $+0.810$, Pearson $+0.926$ on the unseen
points. \emph{Absolute calibration does not}: only $8$ of $52$ fall inside the old
$\pm 0.13$ band, with a systematic per-world offset of $+0.12$ to $+0.15$---the
old line \emph{under}-predicts accuracy, which is at least the safe direction. A
refit pooled over all nine worlds ($n=104$) gives
$\text{accuracy} \approx 1.76 \cdot \mathrm{PX} + 0.13$ with residual $\pm 0.16$.

The obvious suspect for the offset was a confound in our own procedure: the later
worlds switched conditional-independence test, so the shift might have been ours
rather than the world's. Rerunning R1 identically under the new test settles it.
PX and true accuracy move together (mean $+0.046$ and $+0.049$), which slides
points \emph{along} the calibration line rather than above it; the residual goes
from $-0.048$ to $-0.067$, against the $+0.136$ the mechanism worlds show. The
test choice contributes essentially nothing. The offset is a property of the
data-generating domain, which strengthens rather than weakens the conclusion: what
must be recalibrated per domain is the intercept, not the instrument.

The resulting claim, stated at the strength the evidence supports: \textbf{PX is a
rank-reliable selector (out-of-sample Spearman $+0.81$) whose absolute
accuracy calibration is per-domain, to about $\pm 0.16$ F1}. It is not a universal
constant, and quoting it as one would be the easiest way to oversell this work.

Three preconditions emerged from the rungs, and all three are detectable without
ground truth---which is what makes them usable rather than merely known. The
candidate pool must be resolution-homogeneous (the R4 lesson: crossing graphs of
different sizes muddies the pairwise F1, and pool resolution spread is directly
measurable). The dyn channel must be live (an atmosphere-regime rung produced a
dead poke channel, which makes the selector inapplicable rather than wrong; the
number of detected dynamical edges reveals this immediately). And the input series
must be deseasonalized.

\subsection{A by-product: disagreement localizes model deficiencies}

The two channels were built to be crossed, but their \emph{disagreements} turn out
to be informative on their own. On maps whose poke channel is live, the true edges
that appear in the read graph but not in the poke graph are almost exactly the
couplings the forecaster failed to internalize: three specific edges, $15/16$
detections, with no non-ground-truth disagreements. Pooled over all good maps
including those with handicapped poke channels, $62\%$ of disagreements localize a
genuine model deficiency against a $25\%$ base rate.

``In the graph you read but not the graph you poke'' is thus a candidate
model-evaluation product in its own right: it points at where an emulator has
learned the correlational structure of a coupling without implementing the
mechanism. It survives losing the oracle map, so it is available in the GraphCast
setting, where it would name the atmospheric couplings the model predicts through
but does not implement.
